\section{Introduction}
\label{sec:introduction}

The source studied here records multiplication as order: positive integers
are ordered by divisibility, the bottom covers are the active atoms, and
least common multiples are joins.  Incidence inversion then supplies
canonical interval coefficients in the sense of \citet{Rota1964}.  Earlier
steps in the construction showed that a quadratic mixed Gram phase is
analytic but generic, and that local natural counterterms do not make it
arithmetic-selective.  A narrow question remained.  Could a nonlocal
statistic of many atoms, their joint interval, or the entire filtered source
separate integer divisibility from non-arithmetic controls?

The most attractive candidate combines two familiar connectedness
machines.  A tuple of atoms is accepted when its subset joins form a Boolean
interval; the top M{\"o}bius value normalizes the accepted tuple to one.  One
may then apply the partition-lattice cumulant transform of
\citet{Speed1983}.  This construction is label-blind, compatible with
transport, and can inspect pairs or triples without calling a primality
routine.  Yet its two versions fail in opposite directions.  The Boolean
weight is one on both integer prime tuples and formal-generator tuples.  The
connected cumulant of the resulting factorizing moment is zero on the
integer baseline itself.

Those failures are symptoms of a complete obstruction.  The positive
integers under multiplication form a free commutative monoid on their prime
covers.  The valuation vector
\(
  n\mapsto(v_p(n))_p
\)
preserves the bottom, covers, divisibility, lcms, and every interval.  If the
cutoff, roof, Gram, and compiled data visible to the construction are
transported along this map, the integer source and a formal
free-commutative/UFD clone are isomorphic as decorated sources.  An
isomorphism-natural invariant, no matter how global, must take equal values
on them.  Since the desired control statement quantifies over \emph{every}
UFD control, it includes this clone and contradicts baseline nonvanishing.

\begin{figure}[t]
\centering
\begin{tikzpicture}[
  >=Latex,
  source/.style={draw=formalblue,very thick,rounded corners=2pt,fill=formalblue!7,
    minimum height=15mm,text width=45mm,align=center,inner sep=5pt},
  value/.style={draw=deepgreen,thick,rounded corners=2pt,fill=deepgreen!7,
    minimum height=11mm,text width=37mm,align=center},
  bad/.style={draw=stopred,thick,dashed,rounded corners=2pt,fill=stopred!4,
    text width=93mm,align=center,inner sep=5pt},
  every node/.style={font=\small}
]
\node[source] (ints) {decorated integer source\\
  $(\N_{>0},\mid,1,\lcmop;F_N,w,G,\ldots)$};
\node[source,right=18mm of ints] (formal) {transported formal UFD clone\\
  $(\bigoplus_{p\in P}\N e_p,\le,0,\max;F'_N,w',G',\ldots)$};
\draw[<->,very thick,formalblue] (ints) -- node[above,fill=white,inner sep=2pt]
  {$\PhiVal(n)=(v_p(n))_p$} (formal);

\node[value,below=13mm of ints] (iv) {$I(\MZ)$};
\node[value,below=13mm of formal] (fv) {$\PhiVal^* I(\Ffree)$};
\draw[->,thick] (ints) -- node[left] {natural $I$} (iv);
\draw[->,thick] (formal) -- node[right] {natural $I$} (fv);
\draw[<->,double,thick,deepgreen] (iv) -- node[above,fill=white] {equal} (fv);

\node[bad,below=12mm of $(iv)!0.5!(fv)$] (gate)
  {Impossible universal separator:\quad $I(\MZ)\ne0$ but
   $I(U)=0$ for every free-commutative/UFD control $U$};
\draw[stopred,very thick] ($(gate.west)+(4mm,-4mm)$) -- ($(gate.east)+(-4mm,4mm)$);
\end{tikzpicture}

\caption{The valuation map is an isomorphism after every admitted decoration
is transported.  Naturality forces equal outputs.  The barred statement is
therefore an inconsistent universal-control specification, not a challenge
that a more elaborate coefficient can solve.}
\label{fig:clone-obstruction}
\end{figure}

The paper makes four contributions.  First, it states the admissible
decorated source category and proves the free-monoid indistinguishability
theorem without a locality, linearity, additivity, or fixed-arity
hypothesis.  Second, it gives the exact canonicity boundary for the Boolean
weight and proves factorization cancellation for its connected cumulant.
Third, it reports that a strengthened rank-three filter separates all four
finite non-UFD fixtures, then shows why its exact free/UFD clone collision is
decisive.  It also embeds the pair weight into the inherited mixed Gram
series, proves normal convergence on its natural strip, and assigns its
marker ledger to a newly declared functional.  A separate trace-class Gram
matrix owns an auxiliary ordinary determinant without upgrading the chiral
object.  Fourth, the paper converts the obstruction into a branch
closure rule: changing coefficients, cumulant conventions, tuple arity,
cutoffs, or finite parts cannot reopen the same multiplicative source.

The scope is equally important.  The theorem does not forbid an enrichment
that adds independently motivated nonmultiplicative data.  Addition--
multiplication interaction, congruence correspondences, or a genuine
symbolic-dynamical transfer operator would define a new source object and
require a new proof.  The theorem also makes no statement about target
zeros.  Route B remains locked throughout.
